% One page per app; \input by fujinet-palm-guide.tex.

% ---------------------------------------------------------------------------
\newpage
\section{FN HotSync}\label{app:fnhotsync}
\appbanner{FN HotSync}{Browse FujiNet's host slots and queue Palm apps for the next HotSync}
  {\fnbadge\quad creator \texttt{ADHS}\quad FujiNet Classic}

FN HotSync is the Palm's own window on FujiNet Classic's host slots: the
same eight slots an Atari or Apple sees in CONFIG, one of them FujiNet's SD
card and the rest TNFS servers. Browse a server, pick a Palm app or database,
and FN HotSync has FujiNet copy it into \texttt{/palm/install/} on the SD card.
The next HotSync installs it.

\begin{center}
\palmshot{fn-hotsync-main}{The host slots}\hfill
\palmshot{fn-hotsync-browse}{Browsing \texttt{/PALM/} on a TNFS server}\hfill
\palmshot{fn-hotsync-queue}{Queue a file for HotSync}
\end{center}

\begin{multicols}{2}
\subsection*{Using it}
\begin{itemize}
  \item Pick a host slot and tap \btn{Browse}. Folders end in \texttt{/}: select
    one and tap \btn{Open} to go in, \btn{Up} to come back out.
  \item Select a \texttt{.prc}, \texttt{.pdb} or \texttt{.pqa} file and tap
    \btn{Open}; confirm with \btn{Queue}. FujiNet copies it to the SD card.
  \item \btn{Queue} on the main screen lists what is waiting;
    \btn{HotSync Now} there starts the sync.
  \item \btn{Edit} names a slot: \texttt{SD} for FujiNet's card, or a TNFS
    server such as \texttt{apps.irata.online}. FujiNet saves the slots.
  \item \btn{Info} shows FujiNet's WiFi network, IP address and firmware
    version.
  \item \btn{HotSync} closes the cradle port and opens the Palm's HotSync
    app, ready to sync.
\end{itemize}

\subsection*{Good to know}
\begin{itemize}
  \item Queued files go to the slot named \texttt{SD}; without one, FN
    HotSync says so.
  \item The copy runs on FujiNet, server to SD card, so even a large app
    takes only a moment.
  \item A queued file stays until a HotSync installs it; the Palm cannot take
    it back out of the queue.
  \item The NIO version of this app is FI Config, on the next page.
\end{itemize}
\end{multicols}

\begin{center}
\palmshot{fn-hotsync-edit}{Naming a host slot}\hfill
\palmshot{fn-hotsync-queued}{What is waiting for the next sync}\hfill
\palmshot{fn-hotsync-info}{FujiNet's network}
\end{center}

% ---------------------------------------------------------------------------
\newpage
\section{FI Config}\label{app:ficonfig}
\appbanner{FI Config}{Browse TNFS servers and queue Palm apps for HotSync through FujiNet NIO}
  {\fibadge\quad creator \texttt{ANHS}\quad FujiNet NIO}

FI Config is FN HotSync rebuilt for FujiNet NIO. NIO has no host slots, so
FI Config keeps its own list of eight TNFS servers, starting with
\texttt{fujinet.diller.org} and \texttt{apps.irata.online}. Browse one, pick
an app, and FI Config copies it through NIO's file service into
\texttt{hotsync/install/}, ready for the next HotSync.

\begin{center}
\palmshot{fi-config-main}{The TNFS servers}\hfill
\palmshot{fi-config-browse}{Browsing a server}\hfill
\palmshot{fi-config-progress}{Copying into the queue}
\end{center}

\begin{multicols}{2}
\subsection*{Using it}
\begin{itemize}
  \item Pick a server and tap \btn{Browse}; \btn{Open} goes into a folder or
    picks a file, \btn{Up} goes back.
  \item Opening a \texttt{.prc}, \texttt{.pdb} or \texttt{.pqa} asks whether
    to queue it. The copy shows its progress on the line above the list.
  \item \btn{Queue} lists the files waiting in NIO's install folder.
  \item \btn{Edit} changes a server's name; clear it to empty the slot. The
    list lives in the app's own preferences.
  \item \btn{Info} shows NIO's WiFi network, IP address, firmware and board,
    and how many replies were asked for again.
  \item \btn{HotSync} hands the cradle back and opens the HotSync app.
\end{itemize}

\subsection*{Good to know}
\begin{itemize}
  \item The Palm does the copying, a few hundred bytes at a time, so a 30~KB
    app takes about half a minute over TNFS.
  \item If a copy is cut short, HotSync refuses the half-written file rather
    than installing a broken app: queue it again to replace it.
  \item A reply lost on the cable is simply asked for again; Info counts how
    often that happened.
  \item On NIO's console, \texttt{hotsync.fetch} queues an app over HTTP
    without the Palm.
\end{itemize}
\end{multicols}

\begin{center}
\palmshot{fi-config-queued}{The install queue}\hfill
\palmshot{fi-config-info}{NIO's firmware and network}\hfill
\palmshot{fi-hotsync-log}{The HotSync log: installed}
\end{center}

% ---------------------------------------------------------------------------
\newpage
\section{FN Test and FI Test}\label{app:test}
\appbanner{FN Test \textbullet\ FI Test}{Is FujiNet there? A one-tap check of the cradle link}
  {\fnbadge\ \texttt{ADTS}\qquad \fibadge\ \texttt{ANTS}}

The first app to run on a new setup. \btn{Ping} sends FujiNet one request
and shows what comes back. If it says \textsf{FujiNet found!}, the cradle,
the cable, the null-modem adapter and the firmware are all doing their job,
and every other app will work too.

\begin{center}
\palmshot{fn-test}{FN Test finds FujiNet Classic}\hfill
\palmshot{fi-test-ping}{FI Test finds FujiNet NIO}\hfill
\palmshot{fi-test-fetch}{FI Test fetches a web page}
\end{center}

\begin{multicols}{2}
\subsection*{FN Test}
\btn{Ping FujiNet} asks FujiNet Classic for its adapter configuration and
shows the firmware version, FujiNet's IP address and the WiFi network it is
on.

\subsection*{FI Test}
\begin{itemize}
  \item \btn{Ping} asks NIO for its version and board (FujiDevice), then
    for its WiFi status and settings.
  \item \btn{Fetch} gets \texttt{http://example.com/} through NIO's network
    service and shows the HTTP status, the size and the page's title: proof
    that the Palm can reach the internet.
\end{itemize}

\columnbreak
\subsection*{If it fails}
\textsf{No answer from FujiNet} means nothing came back after several tries.
Check the cable and adapter, and that you are running the right app for your
FujiNet: FN Test for Classic, FI Test for NIO. Chapter~\ref{ch:diff} explains
the difference.

\begin{tip}
Right after a HotSync, or after a long pause, the first ping can take a
second longer: FujiNet may be listening for a HotSync at that moment, and the
app simply asks again.
\end{tip}
\end{multicols}

% ---------------------------------------------------------------------------
\newpage
\section{Bounce World}\label{app:bounce}
\appbanner{FN Bounce World \textbullet\ FI Bounce}{Shapes bouncing across a world shared by every connected machine}
  {\fnbadge\ \texttt{ADBW}\qquad \fibadge\ \texttt{ANBW}}

A client for Mark Fisher's \textsf{bounce-world} server. Every machine that
connects (an Atari, a C64, a Palm) gets a window onto one shared world, and
the shapes in it bounce from screen to screen. The Palm draws its part of the
world in block graphics, polling the server several times a second, which
also makes it a fine workout for the cradle link.

\begin{center}
\palmshot{fn-bounce-connect}{Connecting to a server}\hfill
\palmshot{fi-bounce-world}{The world in motion}\hfill
\palmshot{fi-bounce-info}{FI Bounce's link counters}
\end{center}

\begin{multicols}{2}
\subsection*{Using it}
Enter the \textsf{Server} and \textsf{Port} (9003 for the framed protocol of
bounce-world 2.3 and later) and a \textsf{Name} of up to 8 letters, then tap
\btn{Connect}. FI
Bounce starts with \texttt{bouncy.diller.org}; FN Bounce World with
\texttt{localhost}, for FujiNet-PC.

\begin{tabularx}{\linewidth}{@{}l X@{}}
\btn{Frz} \key{F} & freeze or unfreeze the world, for everyone \\
\btn{Rst} \key{R} & reset it \\
\btn{+} \btn{-} & faster, slower (also the scroll keys) \\
\btn{Add} \key{1}--\key{5} & add a body of that size \\
\btn{Who} \key{W} & who else is connected \\
\btn{Info} \key{I} & your name, clients, world size and bodies \\
\key{Q} & disconnect \\
\end{tabularx}

While you are in the world the app buttons play too: 1 freezes, 2 adds a
body, 3 shows who, 4 shows info.

\subsection*{FI Bounce's counters}
FI Bounce's Info adds a line of link statistics: requests, exchanges with
NIO, retries and replies lost. A long run with no losses is a healthy cable.

\begin{warn}[Known issue]
The world keeps drawing while the \menu{World} menu is open, so the menu
shows as a blank box. Its items still work, but the buttons and Graffiti keys
are easier.
\end{warn}
\end{multicols}

% ---------------------------------------------------------------------------
\newpage
\section{Lobby}\label{app:lobby}
\appbanner{FN Lobby \textbullet\ FI Lobby}{Every game server in the FujiNet Game Lobby, in your pocket}
  {\fnbadge\ \texttt{ADLB}\qquad \fibadge\ \texttt{ANLB}}

The FujiNet Game Lobby lists the multiplayer games FujiNet machines can
join: Battleship, 5 Card Stud, Texas Hold'em and many more, with who is
playing and from which platforms. The Lobby app reads the whole list,
\texttt{/viewFull} from \texttt{lobby.fujinet.online}, and shows it on the
Palm.

\begin{center}
\palmshot{fi-lobby-loading}{Reading the lobby}\hfill
\palmshot{fn-lobby-list}{Every server, by game}\hfill
\palmshot{fn-lobby-detail}{One server in detail}
\end{center}

\begin{multicols}{2}
\subsection*{Using it}
\begin{itemize}
  \item The list loads when the app starts. Each row is a game, a server
    and its players, such as \textsf{0/4}.
  \item Tap a row for the server's status, region, players, app key and URL,
    and the platforms it serves with the client URL for each. The scroll keys
    page through it.
  \item \btn{Refresh} reads the list again.
  \item \menu{Lobby > Use QA Lobby} switches to the test lobby, for
    developers; choose it again to go back.
\end{itemize}

\subsection*{Good to know}
\begin{itemize}
  \item The reply is about 30~KB of JSON, more than a Palm would like to hold
    at once. The app reads it as it streams in, keeping only what it shows,
    up to 48 servers.
  \item FI Lobby stays usable while the list loads: the menus, Refresh and
    Details all answer straight away.
\end{itemize}
\end{multicols}

% ---------------------------------------------------------------------------
\newpage
\section{Texas Hold'em}\label{app:holdem}
\appbanner{FN Texas Hold'em \textbullet\ FI Texas Hold'em}{Multiplayer poker, at the same tables as Atari, Apple II and CoCo players}
  {\fnbadge\ \texttt{ADTH}\qquad \fibadge\ \texttt{ANTH}}

A Palm client for the FujiNet Texas Hold'em server, listed in the Game Lobby.
Sit at a table with AI players or with friends on other machines; the server
deals and runs every rule, and the Palm shows the felt, the cards and the
moves you can make.

\begin{center}
\palmshot{fn-holdem-setup}{Pick a name and a table}\hfill
\palmshot{fn-holdem-table}{Your move}\hfill
\palmshot{fi-holdem-showdown}{The showdown}
\end{center}

\begin{multicols}{2}
\subsection*{Using it}
\begin{itemize}
  \item Enter a \textsf{Name} (up to 8 letters and digits). The tables on the
    \textsf{Server} are listed with their players; pick one (or type its name
    under \textsf{Table}) and tap \btn{Sit Down}.
  \item Your two cards sit beside the board; the other players are above, face
    down until the showdown.
  \item On your turn the moves appear along the bottom with the server's
    clock. Tap one, or pick one with the scroll keys and press any app button.
  \item Graffiti works too: \key{f} fold, \key{c} check or call, \key{b} bet,
    \key{r} raise, \key{a} all-in.
  \item Tap the status line, or \menu{Table > Last Result}, for the last
    hand. \menu{Table > Leave Table} gives up your seat.
\end{itemize}

\subsection*{Good to know}
\begin{itemize}
  \item The default server is \texttt{th.carr-designs.com}. A lobby address
    that ends in \texttt{?table=den} can be typed as the server, and the table
    is taken from it.
  \item The cards are the Apple~II client's art, which needs no colour; the
    small cards are the Atari Lynx client's.
  \item Leaving the app, or switching to another, leaves the table for you.
\end{itemize}
\end{multicols}

% ---------------------------------------------------------------------------
\newpage
\section{FujiPet}\label{app:fujipet}
\appbanner{FN FujiPet \textbullet\ FI FujiPet}{A virtual pet that grows up like the 1996 original, and makes friends online}
  {\fnbadge\ \texttt{ADFP}\qquad \fibadge\ \texttt{ANFP}}

Name your pet and pick a Cat, Dog, Bunny or Dino; the egg hatches five
minutes later. It grows from baby to child, teen and adult over a few days,
and how well it is looked after decides what it grows into and how long it
lives. With \textsf{Meet friends online} on, it keeps a connection to a
FujiPet friend server and can chat with pets on other Palms.

\begin{center}
\palmshot{fn-fujipet-new}{A new egg}\hfill
\palmshot{fn-fujipet-main}{All grown up, in colour}\hfill
\palmshot{fn-fujipet-friends}{Friends online}
\end{center}

\begin{multicols}{2}
\subsection*{Looking after it}
The eight icons are the original's. Along the top: \textbf{Feed} (a meal
fills a hunger heart, a snack a happy one), \textbf{Light}, \textbf{Play}
(guess which way it will look, three of five to win) and \textbf{Medicine}.
Along the bottom: \textbf{Clean}, \textbf{Status}, \textbf{Scold} and
\textbf{Friends}. Tap the pet itself to make it happy. A heart running empty makes it call
(the \textsf{!} in the corner); a call left for 15 minutes is a care mistake.
Mess and too many sweets make it sick. The pet lives on while the app is
closed, and catches up when you return.

\subsection*{Friends}
\btn{Register} adds your pet to the friend server. The Friends screen lists
the pets there, online ones first in bold; pick one, choose one of eight
messages under \textsf{Say} and tap \btn{Send}. It pops up as a bubble on
their Palm at once, and a kind one cheers their pet up.

\subsection*{Settings}
\menu{Options > Settings} sets the friend server (\textsf{host:port}), turns
going online and sounds on or off, and runs time \textsf{Fast} (a pet minute
a second) or \textsf{Turbo} (a pet day a minute) to watch it grow up.
\end{multicols}

\begin{center}
\palmshot{fn-fujipet-status}{Status: hearts, age and weight}\hfill
\palmshot{fi-fujipet-message}{A message from a friend}\hfill
\palmshot{fn-fujipet-mono}{In greys on a Palm V}
\end{center}
